\documentclass[11pt, a4paper]{article}

\usepackage[a4paper, top=2.5cm, bottom=2.5cm, left=2cm, right=2cm]{geometry}
\usepackage{amsmath}
\usepackage{amssymb}
\usepackage[colorlinks=true, linkcolor=blue, urlcolor=blue, citecolor=blue]{hyperref}

\title{Theory of Everything - Volume 44 \\ \Large Universal Expansion \& The Big Rip}
\author{Omega Protocol}
\date{\today}

\begin{document}

\maketitle

\section*{Layman's Summary}
How will the universe end? Volume 44 looks at the far, far future. We prove that while the universe will continue to expand and cool down into total darkness, the actual fabric of reality will never tear apart, and incredibly efficient thought and computation can survive forever.

\section{The Thermodynamic End State (Axiom 47)}
We establish that the continuous, accelerating expansion of the universe is simply reality coasting toward its final state of maximum entropy—the "Heat Death." Eventually, all stars will burn out, all black holes will evaporate, and the universe will be a cold, dark, perfectly balanced soup of radiation.

\section{The Big Rip Avoidance (Theorem)}
Some physicists fear that dark energy will eventually grow so strong it will rip atoms and spacetime itself apart (the "Big Rip"). We prove this is impossible. The fundamental "informational stiffness" of the universe's algebra acts like an unbreakable rubber band. It can stretch forever, but it cannot mathematically snap.

\section{Late-stage Computation (Theorem)}
Can life survive the Heat Death? "Landauer's Principle" states that the colder the environment, the less energy it takes to think. We prove that as the universe freezes, the computational efficiency of any surviving digital intelligence will increase exponentially. They will think slower, but they can think forever on almost zero energy.

\end{document}
